\centerline{\Huge \bf  Олимпиада II}
\centerline{\Huge \bf 8 класс}

\begin{flushright}
    \scriptsize{Заключительный отбор в Сириус 2022г.\\ Кто загуглит - тому бан!}
\end{flushright}

\problem{1}
Никита загадал четырёхзначное число. Известно, что цифра в разряде сотен меньше одной из цифр этого числа на $5$, больше другой на $4$, а от цифры в разряде десятков отличается на $2$. Какое число мог загадать Никита? Не забудьте указать все варианты и объяснить, что других нет.

\problem{2}
Пусть $x$ и $y$ $-$ положительные действительные числа, для которых $x^2 +y^2 = \dfrac{1}{2}$. Докажите, что 
\[\dfrac{1}{1+2x} + \dfrac{1}{1+2y} \geq 1\]
При каких условиях неравенство обращается в равенство?

\problem{3}
В треугольнике ABC на отрезке BC отмечены точки $D$ и $E$ такие, что $BD =
CE \ \text{и} \ \\ \angle BAD =  \angle CAE$. Докажите, что треугольник $ABC$ равнобедренный.

\problem{4}
В квадрате $45 \times 45$ покрашено $2022$ клетки. Требуется выбрать квадрат $n \times n$ в котором все клетки покрашены. Найдите наибольшее $n$ при котором можно выбрать указанный квадрат для любого расположения покрашенных клеток.

\problem{5}
Олеся нарисовала пятиугольник. Петя измерил длины его сторон и получил $5$ целых чисел, сумма которых нечётна. Какое наибольшее количество прямых углов могло быть в таком пятиугольнике?

\problem{6}
Антон выбирает в качестве начального числа целое число $n \geq 0$, которое не является квадратом. Берта прибавляет к этому числу $n + 1.$ Если эта сумма является полным квадратом, Берта выиграла. В противном случае Антон прибавляет к полученной сумме, последующее число $n + 2$. Если эта сумма является полным квадратом, он выиграл. В противном случае $-$ снова ход Берты, и она прибавляет следующее число $n + 3$, и так далее. Докажите, что существует бесконечное множество чисел $n$, начиная с которых Антон победит